\setcounter{section}{16}

  \zwischenueberschrift{Übungsaufgaben}

\inputaufgabegibtloesung
{}
{

Wir betrachten die Standardparabel

\mavergleichskette
{\vergleichskette
{ Y
}
{ \subseteq }{ \R^2
}
{  }{
}
{    }{
}
{   }{
}
}
{}{}{}
mit der induzierten riemannschen Struktur und mit der Parametrisierung
\maabbeledisp {\varphi} {\R} { Y
} { t } { \left(  t   , \,  t^2                   \right)
} {,}
die wir als eine
\zusatzklammer {inverse} {} {}
Karte betrachten. Bestimme die riemannsche
\definitionsverweis {Fundamentalfunktion}{}{}

\mavergleichskette
{\vergleichskette
{ g(t)
}
{ = }{ g_{11} (t)
}
{  }{
}
{    }{
}
{   }{
}
}
{}{}{.}

}
{} {}

\inputaufgabe
{}
{

Es sei $M$ eine
\definitionsverweis {differenzierbare Mannigfaltigkeit}{}{}
derart, dass das
\definitionsverweis {Tangentialbündel}{}{}
$TM$ trivial ist. Zeige, dass man über eine differenzierbare Trivialisierung

\mavergleichskette
{\vergleichskette
{ TM
}
{ \cong }{ M \times \R^n
}
{  }{
}
{    }{
}
{   }{
}
}
{}{}{}
\zusatzklammer {mithilfe des
\definitionsverweis {Standardskalarproduktes}{}{}
auf dem $\R^n$} {} {}
$M$ zu einer
\definitionsverweis {riemannschen Mannigfaltigkeit}{}{}
machen kann.

}
{} {}

\inputaufgabe
{}
{

Wir betrachten den Kreis

\mavergleichskette
{\vergleichskette
{ S^1
}
{ \subseteq }{ \R^2
}
{  }{
}
{    }{
}
{   }{
}
}
{}{}{}
mit der
\definitionsverweis {induzierten riemannschen Struktur}{}{.}
Die Trivialisierung des
\definitionsverweis {Tangentialbündels}{}{}
\maabbeledisp {} { S^1 \times \R } { T S^1
} {( a,b, u )} { (a,b, -ub,ua )
} {,}
führt über das Standardskalarprodukt von $\R$  ebenfalls zu einer riemannschen Struktur auf $S^1$. Zeige, dass beide riemannschen Strukturen übereinstimmen.

}
{} {}

\inputaufgabe
{}
{

Es seien
\mathkor {} {L} {und} {M} {}
\definitionsverweis {riemannsche Mannigfaltigkeiten}{}{.}
Zeige, dass $L \times M$ in natürlicher Weise ebenfalls eine riemannsche Mannigfaltigkeit ist.

}
{} {}

\inputaufgabe
{}
{

Wir betrachten eine offene Menge
\mathl{V \subseteq \R^n}{} als
\definitionsverweis {riemannsche Mannigfaltigkeit}{}{.} Was ist die
\definitionsverweis {kanonische Volumenform}{}{} auf $V$?

}
{} {}

\inputaufgabe
{}
{

Es sei $M$ eine
\definitionsverweis {riemannsche Mannigfaltigkeit}{}{.}
Zeige, dass die
\definitionsverweis {Abbildung}{}{}
\maabbeledisp {} {
{ \mathcal V } ( M )  } {
{ \mathcal E }^{  1 } (  M   )
} { F } { \omega_F
} {,}
mit

\mavergleichskettedisp
{\vergleichskette
{ \left(  \omega_F(P)  \right) (v)
}
{ \defeq} { \left\langle F(P) , v  \right\rangle_P
}
{ } {
}
{ } {
}
{ } {
}
}
{}{}{,}
\definitionsverweis {linear}{}{}
ist.

}
{} {}

\inputaufgabe
{}
{

Wir betrachten eine offene Menge
\mathl{V \subseteq \R^n}{} als
\definitionsverweis {riemannsche Mannigfaltigkeit}{}{.} Was besagt die in
[[Riemannsche Mannigfaltigkeit/Vektorfelder und 1-Formen/Fakt|Lemma 16.3]] beschriebene Korrespondenz zwischen
\definitionsverweis {Vektorfeldern}{}{}
und
$1$-\definitionsverweis {Differentialformen}{}{} in dieser Situation?

}
{} {}

\inputaufgabe
{}
{

Begründe
[[Abgeschlossene Untermannigfaltigkeit/R^n/Einschränkung eines Vektorfeldes/Bemerkung|Bemerkung 16.4]].

}
{} {}

\inputaufgabe
{}
{

Es sei $M$ eine
\definitionsverweis {orientierte}{}{}
\definitionsverweis {riemannsche Mannigfaltigkeit}{}{.} Zeige, dass die
\definitionsverweis {kanonische Volumenform}{}{} $\omega$ dadurch festgelegt ist, dass sie in jedem Punkt für eine die Orientierung repräsentierende Orthonormalbasis den Wert $1$ besitzt.

}
{} {}

\inputaufgabe
{}
{

Es sei $V$ ein $n$-dimensionaler reeller orientierter Vektorraum und $\lambda$ ein
\definitionsverweis {translationsinvariantes Maß}{}{}
auf $V$. Zeige, dass die Zuordnung
\maabbeledisp {} {V \times \cdots \times V} {\R
} {(v_1 , \ldots , v_n) } { \pm \lambda (P(v_1 , \ldots , v_n))
} {,}
wobei das Vorzeichen positiv zu wählen ist, wenn die Vektoren die Orientierung repräsentieren, eine alternierende multilineare Abbildung ist.

}
{} {}

\inputaufgabe
{}
{

Zeige, dass bei einer
\definitionsverweis {riemannschen Mannigfaltigkeit}{}{}
die
\definitionsverweis {Kartenabbildungen}{}{}
\maabbdisp {\alpha} { U } { V
} {}
im Allgemeinen keine
\definitionsverweis {Isometrie}{}{}
\maabbdisp {T_P(\alpha)} {T_PU} {T_{\alpha(P)} V
} {}
induzieren
\zusatzklammer {wenn
\mathl{T_PU}{} mit
\mathl{\left\langle  - ,  -  \right\rangle_P}{} und
\mathl{T_{\alpha(P)} V= \R^n}{} mit dem Standardskalarprodukt versehen ist} {} {.}

}
{} {}

\inputaufgabegibtloesung
{}
{

Wir betrachten den Graph $M$ der Abbildung
\maabbeledisp {\varphi} {\R^2} {\R
} {(u,v)} { u^2+uv-v^3
} {,} als zweidimensionale abgeschlossene Untermannigfaltigkeit des $\R^3$, also

\mavergleichskettedisp
{\vergleichskette
{ M
}
{ =} { \left\{ (u,v,u^2+uv-v^3)  \mid  (u,v) \in \R^2         \right\}
}
{ } {
}
{ } {
}
{ } {
}
}
{}{}{}
mit der vom $\R^3$ induzierten riemannschen Metrik. Es sei
\maabbeledisp {\psi} {\R^2} { M
} {(u,v)} { (u,v,u^2+uv-v^3)
} {,}
die zugehörige Diffeomorphie.

a) Bestimme das totale Differential zu $\psi$ sowie die Bildvektoren
\mathl{T_P(\psi) (e_1)}{} und
\mathl{T_P(\psi) (e_2)}{} in
\mathl{T_{\psi(P)}M}{.}

b) Bestimme für jeden Punkt der Form

\mavergleichskette
{\vergleichskette
{ P
}
{ = }{ (u,0)
}
{  }{
}
{    }{
}
{   }{
}
}
{}{}{}
den Flächeninhalt des von
\mathl{T_P(\psi) (e_1)}{} und
\mathl{T_P(\psi) (e_2)}{} in
\mathl{T_{\psi(P)}M}{} aufgespannten Parallelogramms.

c) Bestimme für jeden Punkt der Form

\mavergleichskette
{\vergleichskette
{ P
}
{ = }{ (0,v)
}
{  }{
}
{    }{
}
{   }{
}
}
{}{}{}
den Flächeninhalt des von
\mathl{T_P(\psi) (e_1)}{} und
\mathl{T_P(\psi) (e_2)}{} in
\mathl{T_{\psi(P)}M}{} aufgespannten Parallelogramms.

}
{} {}

\inputaufgabe
{}
{

Es sei
\maabbdisp {\varphi} {\R^n } {\R
} {}
\zusatzklammer {mit \mathlk{m=n -1 \geq 0}{}} {} {}
eine
\definitionsverweis {stetig differenzierbare Funktion}{}{,}
die in jedem Punkt der
\definitionsverweis {Faser}{}{}
$M$ über
\mathl{0 \in \R}{}
\definitionsverweis {regulär}{}{}
sei. Wir fassen $M$ als eine orientierte riemannsche Mannigfaltigkeit auf. Zeige, dass zwischen der Volumenform $\tau$ aus
[[Differenzierbare reguläre Abbildung/R^n/Faser besitzt Volumenform über Gradienten/Fakt|Korollar 15.6]]
und der
\definitionsverweis {kanonischen Volumenform}{}{}
$\omega$ die Beziehung

\mavergleichskettedisp
{\vergleichskette
{\tau(P,v_1 , \ldots , v_m)
}
{ =} { \pm \Vert {
\operatorname{Grad} \, \varphi ( P ) } \Vert  \omega(P, v_1 , \ldots , v_m)
}
{ } {
}
{ } {
}
{ } {
}
}
{}{}{}
besteht.

}
{} {}

\inputaufgabe
{}
{

Es sei
\maabbdisp {\varphi} {\R^n } {\R^\ell
} {}
\zusatzklammer {mit \mathlk{m=n - \ell \geq 0}{}} {} {}
eine
\definitionsverweis {stetig differenzierbare Abbildung}{}{,}
die in jedem Punkt der
\definitionsverweis {Faser}{}{}
$M$ über
\mathl{0 \in \R^\ell}{}
\definitionsverweis {regulär}{}{}
sei. Wir fassen $M$ als eine orientierte riemannsche Mannigfaltigkeit auf. Es sei vorausgesetzt, dass die Gradienten

\mathdisp {\operatorname{Grad} \, \varphi_1  (  P ) , \ldots ,
\operatorname{Grad} \, \varphi_\ell ( P )} {  }

für jeden Punkt von
\mathl{P\in M}{} senkrecht aufeinander stehen. Zeige, dass zwischen der Volumenform $\tau$ aus
[[Differenzierbare reguläre Abbildung/R^n/Faser besitzt Volumenform über Gradienten/Fakt|Korollar 15.6]]
und der
\definitionsverweis {kanonischen Volumenform}{}{}
$\omega$ die Beziehung

\mavergleichskettedisp
{\vergleichskette
{\tau(P,v_1 , \ldots , v_m)
}
{ =} { \pm \Vert {
\operatorname{Grad} \, \varphi_1  (  P ) } \Vert  \cdots \Vert {
\operatorname{Grad} \, \varphi_\ell ( P ) } \Vert \omega(P, v_1 , \ldots , v_m)
}
{ } {
}
{ } {
}
{ } {
}
}
{}{}{}
besteht.

}
{} {}

\inputaufgabe
{}
{

Zeige, dass die
\definitionsverweis {kanonische Flächenform}{}{}
auf der
\definitionsverweis {Einheitssphäre}{}{}

\mavergleichskette
{\vergleichskette
{S^2
}
{ \subset }{ \R^3
}
{  }{
}
{    }{
}
{   }{
}
}
{}{}{}
gleich

\mathdisp {x dy \wedge dz  -y dx \wedge dz +  z dx \wedge dy} {  }

ist, wobei $x,y,z$ die Koordinaten des $\R^3$ seien.

}
{} {}

Ein
\definitionsverweis {Vektorbündel}{}{}
\maabb {p} { E } { X
} {}
über einer
\definitionsverweis {differenzierbare Mannigfaltigkeit}{}{}
$X$ heißt
\definitionswort {riemannsches Vektorbündel}{,}
wenn auf jeder Faser $E_x$ ein
\definitionsverweis {Skalarprodukt}{}{}
erklärt ist mit der Eigenschaft, dass es trivialisierende
\definitionsverweis {Karten}{}{}
\maabbdisp {\alpha} { U } { V \subseteq \R^n
} {}
mit
\maabbdisp {\theta} {E\vert_U \cong U \times \R^r } { V \times \R^r
} {}
derart gibt, dass

\mavergleichskettedisp
{\vergleichskette
{ h_{ij} (Q )
}
{ \defeq} { \left\langle  \theta^{-1}(Q,e_i)  , \theta^{-1}(Q,e_j)    \right\rangle_{\alpha^{-1}(Q)}
}
{ } {
}
{ } {
}
{ } {
}
}
{}{}{}
stetig differenzierbar sind.

Bei einer
\definitionsverweis {riemannschen Mannigfaltigkeit}{}{}
ist also das
\definitionsverweis {Tangentialbündel}{}{}
ein riemannsches Vektorbündel.

\inputaufgabe
{}
{

Es sei
\maabb {p} { E } { M
} {}
ein
\definitionsverweis {riemannsches Vektorbündel}{}{}
über einer
\definitionsverweis {differenzierbare Mannigfaltigkeit}{}{}
$M$ und sei
\maabb {\varphi} { L } { M
} {}
eine
\definitionsverweis {differenzierbare Abbildung}{}{}
mit dem
\definitionsverweis {zurückgezogenen Vektorbündel}{}{}
\maabb {} {\varphi^*E} {L
} {.}
Zeige, dass $\varphi^*E$ in natürlicher Weise ebenfalls ein riemannsches Vektorbündel ist.

}
{} {}

  \zwischenueberschrift{Aufgaben zum Abgeben}

Bei einer riemannschen Mannigfaltigkeit $M$ definiert man zu einem Tangentialvektor

\mavergleichskette
{\vergleichskette
{ v
}
{ \in }{ T_PM
}
{  }{
}
{    }{
}
{   }{
}
}
{}{}{}
die Norm durch

\mavergleichskette
{\vergleichskette
{ \Vert {v} \Vert
}
{ = }{ \sqrt{ \left\langle v , v  \right\rangle_P }
}
{  }{
}
{    }{
}
{   }{
}
}
{}{}{.}

\inputaufgabe
{4}
{

Es sei $M$ eine
\definitionsverweis {riemannsche Mannigfaltigkeit}{}{.}
Zeige, dass die Zuordnung
\maabbeledisp {} {TM} {\R
} { v } { \Vert { v } \Vert
} {,}
\definitionsverweis {stetig}{}{}
ist.

}
{} {}

\inputaufgabe
{3}
{

Zeige, dass
\mathl{\R \times \R_+}{} mit der durch die
\definitionsverweis {Hesse-Form}{}{} zur
\definitionsverweis {Funktion}{}{}
\maabbeledisp {f} {\R \times \R_+} {\R
} {(x,y)} { x^2+y^4
} {,}
gegebenen Bilinearform eine
\definitionsverweis {riemannsche Mannigfaltigkeit}{}{}
ist.

}
{} {}

\inputaufgabe
{4}
{

Man gebe für jeden Punkt
\mathl{P=(x,y,z)}{} der
\definitionsverweis {Einheitssphäre}{}{}
$K$ eine
\definitionsverweis {Orthonormalbasis}{}{}
in
\mathl{T_PK \subset \R^3}{} an
\zusatzklammer {bezüglich der induzierten
\definitionsverweis {riemannschen Struktur}{}{}} {} {.}

}
{} {}

\inputaufgabe
{6}
{

Im $\R^3$ sei das
\definitionsverweis {Ellipsoid}{}{}

\mathdisp {E= \left\{ (x,y,z)  \mid   x^2+y^2+3z^2 \leq 5        \right\}} {  }
 und die Ebene

\mathdisp {M= \left\{ (x,y,z)  \mid   7x-3y-2z =  2         \right\}} {  }
 gegeben. Berechne den Flächeninhalt des Durchschnitts
\mathl{M \cap E}{.}

}
{} {}

\inputaufgabe
{6}
{

Man erstelle eine Computergraphik, die die in
[[Abgeschlossene Untermannigfaltigkeit/R^n/Einschränkung eines Vektorfeldes/Bemerkung|Bemerkung 16.4]]
beschriebene Situation anhand einer Fläche im $\R^3$ veranschaulicht.

}
{} {}

 [[Kategorie:Latexseite]]
